% Options for packages loaded elsewhere
\PassOptionsToPackage{unicode}{hyperref}
\PassOptionsToPackage{hyphens}{url}
\documentclass[
]{article}
\usepackage{xcolor}
\usepackage{amsmath,amssymb}
\setcounter{secnumdepth}{-\maxdimen} % remove section numbering
\usepackage{iftex}
\ifPDFTeX
  \usepackage[T1]{fontenc}
  \usepackage[utf8]{inputenc}
  \usepackage{textcomp} % provide euro and other symbols
\else % if luatex or xetex
  \usepackage{unicode-math} % this also loads fontspec
  \defaultfontfeatures{Scale=MatchLowercase}
  \defaultfontfeatures[\rmfamily]{Ligatures=TeX,Scale=1}
\fi
\usepackage{lmodern}
\ifPDFTeX\else
  % xetex/luatex font selection
\fi
% Use upquote if available, for straight quotes in verbatim environments
\IfFileExists{upquote.sty}{\usepackage{upquote}}{}
\IfFileExists{microtype.sty}{% use microtype if available
  \usepackage[]{microtype}
  \UseMicrotypeSet[protrusion]{basicmath} % disable protrusion for tt fonts
}{}
\makeatletter
\@ifundefined{KOMAClassName}{% if non-KOMA class
  \IfFileExists{parskip.sty}{%
    \usepackage{parskip}
  }{% else
    \setlength{\parindent}{0pt}
    \setlength{\parskip}{6pt plus 2pt minus 1pt}}
}{% if KOMA class
  \KOMAoptions{parskip=half}}
\makeatother
\usepackage{longtable,booktabs,array}
\newcounter{none} % for unnumbered tables
\usepackage{calc} % for calculating minipage widths
% Correct order of tables after \paragraph or \subparagraph
\usepackage{etoolbox}
\makeatletter
\patchcmd\longtable{\par}{\if@noskipsec\mbox{}\fi\par}{}{}
\makeatother
% Allow footnotes in longtable head/foot
\IfFileExists{footnotehyper.sty}{\usepackage{footnotehyper}}{\usepackage{footnote}}
\makesavenoteenv{longtable}
\usepackage{graphicx}
\makeatletter
\newsavebox\pandoc@box
\newcommand*\pandocbounded[1]{% scales image to fit in text height/width
  \sbox\pandoc@box{#1}%
  \Gscale@div\@tempa{\textheight}{\dimexpr\ht\pandoc@box+\dp\pandoc@box\relax}%
  \Gscale@div\@tempb{\linewidth}{\wd\pandoc@box}%
  \ifdim\@tempb\p@<\@tempa\p@\let\@tempa\@tempb\fi% select the smaller of both
  \ifdim\@tempa\p@<\p@\scalebox{\@tempa}{\usebox\pandoc@box}%
  \else\usebox{\pandoc@box}%
  \fi%
}
% Set default figure placement to htbp
\def\fps@figure{htbp}
\makeatother
\setlength{\emergencystretch}{3em} % prevent overfull lines
\providecommand{\tightlist}{%
  \setlength{\itemsep}{0pt}\setlength{\parskip}{0pt}}
\usepackage{bookmark}
\IfFileExists{xurl.sty}{\usepackage{xurl}}{} % add URL line breaks if available
\urlstyle{same}
\hypersetup{
  hidelinks,
  pdfcreator={LaTeX via pandoc}}

\author{}
\date{}

\begin{document}

\begin{quote}
PARQUE CIENT˝FICO Y TECNOLÓGICO UTPL
\end{quote}

Análisis de Solución del Reto

\begin{quote}
Plataforma SaaS de Distribución con Vendedores Freelance
\end{quote}

Institución: Parque Científico y Tecnológico UTPL

Versión: 1.0

Fecha: Mayo 2026

Tecnología: Django (Python)

Documento generado en L\textsuperscript{A}TEX PRESENTACION-1BIM

\section{Contents}\label{contents}

Capítulo 1: Introducción y Contexto 2

0.1 Introduccion . . . . . . . . . . . . . . . . . . . . . . . . . . . .
. . . . . . . . . . . . . . . . . . . . . . . . . . 2

0.2 Problemática. . . . . . . . . . . . . . . . . . . . . . . . . . . .
. . . . . . . . . . . . . . . . . . . . . . . . . . 2

0.3 Metodología o Estrategia de Desarrollo . . . . . . . . . . . . . . .
. . . . . . . . . . . . . . . . . 2

Capítulo 2: Planificación y Análisis 4

0.4 Análisis . . . . . . . . . . . . . . . . . . . . . . . . . . . . . .
. . . . . . . . . . . . . . . . . . . . . . . . . . . . . 4

0.5 Requisitos Funcionales . . . . . . . . . . . . . . . . . . . . . . .
. . . . . . . . . . . . . . . . . . . . . . . 4

0.6 Requisitos No Funcionales . . . . . . . . . . . . . . . . . . . . .
. . . . . . . . . . . . . . . . . . . . . . 5

0.7 Historias de Usuario / Casos de Uso . . . . . . . . . . . . . . . .
. . . . . . . . . . . . . . . . . . . 6

0.7.1 Actores del Sistema . . . . . . . . . . . . . . . . . . . . . . .
. . . . . . . . . . . . . . . . . . . 6

0.7.2 Casos de Uso - Módulo de Pedidos. . . . . . . . . . . . . . . . .
. . . . . . . . . . . . . 7

0.7.3 Historias de Usuario Representativas. . . . . . . . . . . . . . .
. . . . . . . . . . . . . . 7

0.8 Sprint. . . . . . . . . . . . . . . . . . . . . . . . . . . . . . .
. . . . . . . . . . . . . . . . . . . . . . . . . . . . . . . 8

0.9 Diagramas: Clases y BD (Diseño Lógico) . . . . . . . . . . . . . . .
. . . . . . . . . . . . . . . . . . . 8

0.9.1 Descripción del Diagrama de Clases . . . . . . . . . . . . . . . .
. . . . . . . . . . . . . . 8

0.9.2 Preguntas Guía para el Diagrama de Clases y BD . . . . . . . . . .
. . . . . . . . . 9

0.10 Diagrama de Componentes. . . . . . . . . . . . . . . . . . . . . .
. . . . . . . . . . . . . . . . . . . . . . 9

0.11 Arquitectura Lógica y Física (Propuesta) . . . . . . . . . . . . .
. . . . . . . . . . . . . 10

0.11.1 Arquitectura Lógica. . . . . . . . . . . . . . . . . . . . . . .
. . . . . . . . . . . . . . . . . . . . 10

0.11.2 Arquitectura Física (Propuesta de Despliegue) . . . . . . . . . .
. . . . . . . . . . 10

0.11.3 Flujo de Proceso del Pedido . . . . . . . . . . . . . . . . . . .
. . . . . . . . . . . . . . . . 11

0.12 Prototipo de Pantallas No Funcional. . . . . . . . . . . . . . . .
. . . . . . . . . . . . . . . . 11

Referencias 12

\section{Capítulo 1: Introducción y
Contexto}\label{capuxedtulo-1-introducciuxf3n-y-contexto}

0.1 Introducción

El presente documento constituye el análisis formal de la solución
propuesta para el reto planteado por el Parque Científico y Tecnológico
de la Universidad Técnica Particular de Loja (UTPL). La solución aborda
una problemática concreta de distribución n comercial en zonas rurales y
semi-urbanas de Latinoamérica, mediante el diseño de una plataforma SaaS
(Software as a Service) que conecta empresas proveedoras, vendedores
freelance y tiendas o negocios destino.

Este informe documenta los requisitos funcionales y no funcionales
identificados, los casos de uso, los diagramas preliminares y la
arquitectura técnica propuesta, con especial énfasis en el uso del
framework Django (Python) como tecnología principal de desarrollo
backend.

0.2 Problemática

La brecha de abastecimiento que enfrentan tiendas y negocios en zonas
rurales o semiurbanas de Latinoamérica constituye la problemática
central del presente reto. Estos establecimientos presentan dificultades
estructurales para acceder a proveedores mayoristas, originadas
principalmente por:

\begin{quote}
Falta de conectividad a internet o inestabilidad de red o de banda
ancha.

Barreras de alfabetización digital entre los propietarios y empleados de
los negocios.

Ausencia de canales de distribución eficientes y formalizados.

Distancias geográficas que encarecen la logística tradicional.
\end{quote}

Como consecuencia, muchos pequeños negocios dependen de intermediarios
informales, pagan precios más altos o simplemente quedan desabastecidos
de productos clave para su operación.

0.3 Metodología o Estrategia de Desarrollo

Para el desarrollo de la solución se adoptará la metodología ágil Scrum,
organizada en sprints quincenales. Esta metodología permite entregas
incrementales y veri cables, facilitando la validación temprana de
funcionalidades con los actores del sistema.

Herramientas y tecnologías de trabajo:

{\def\LTcaptype{none} % do not increment counter
\begin{longtable}[]{@{}
  >{\raggedright\arraybackslash}p{(\linewidth - 2\tabcolsep) * \real{0.2329}}
  >{\raggedright\arraybackslash}p{(\linewidth - 2\tabcolsep) * \real{0.6407}}@{}}
\toprule\noalign{}
\begin{minipage}[b]{\linewidth}\raggedright
\begin{quote}
Componente
\end{quote}
\end{minipage} & \begin{minipage}[b]{\linewidth}\raggedright
Tecnología / Herramienta
\end{minipage} \\
\midrule\noalign{}
\endhead
\bottomrule\noalign{}
\endlastfoot
\begin{minipage}[t]{\linewidth}\raggedright
\begin{quote}
Backend principal
\end{quote}
\end{minipage} & Django 5.x (Python 3.12) framework web de alto nivel \\
\begin{minipage}[t]{\linewidth}\raggedright
\begin{quote}
API REST
\end{quote}
\end{minipage} & Django REST Framework (DRF) serializers, viewsets,
routers \\
\begin{minipage}[t]{\linewidth}\raggedright
\begin{quote}
Base de datos
\end{quote}
\end{minipage} & PostgreSQL 16 motor relacional con soporte ACID
completo \\
\begin{minipage}[t]{\linewidth}\raggedright
\begin{quote}
Autenticación
\end{quote}
\end{minipage} & JWT mediante djangorestframework-simplejwt + 2FA \\
\begin{minipage}[t]{\linewidth}\raggedright
\begin{quote}
Frontend
\end{quote}
\end{minipage} & React 18 + Tailwind CSS SPA con vistas diferenciadas
por rol \\
\begin{minipage}[t]{\linewidth}\raggedright
\begin{quote}
API Gateway
\end{quote}
\end{minipage} & Nginx proxy inverso, rate limiting, terminación TLS \\
\begin{minipage}[t]{\linewidth}\raggedright
\begin{quote}
Pasarela de pagos
\end{quote}
\end{minipage} & Integración con proveedor PCI-DSS (Stripe /
PayPhone) \\
\begin{minipage}[t]{\linewidth}\raggedright
\begin{quote}
Notificaciones
\end{quote}
\end{minipage} & Celery + Redis tareas asíncronas; SMTP / SendGrid para
emails \\
\begin{minipage}[t]{\linewidth}\raggedright
\begin{quote}
Gestión de

proyectos
\end{quote}
\end{minipage} & GitHub Projects + tablero Kanban por sprint \\
\begin{minipage}[t]{\linewidth}\raggedright
\begin{quote}
Prototipado UI
\end{quote}
\end{minipage} & Figma, wireframes y prototipo no funcional de
pantallas \\
\begin{minipage}[t]{\linewidth}\raggedright
\begin{quote}
Control de versiones
\end{quote}
\end{minipage} & Git + GitHub (repositorio del equipo) \\
\begin{minipage}[t]{\linewidth}\raggedright
\begin{quote}
Documentación API
\end{quote}
\end{minipage} & drf-spectacular (OpenAPI 3 / Swagger UI) \\
\end{longtable}
}

Table 1: Stack tecnológico del proyecto

¿Por qué Django?

Django es un framework maduro y con baterías incluidas (batteries
included) que ofrece: ORM potente compatible con PostgreSQL, sistema de
autenticación y permisos integrado, panel de administración automático,
soporte nativo para migraciones de base de datos y una comunidad activa.
Su ecosistema con Django REST Framework permite construir APIs RESTful
documentadas de manera eficiente, cumpliendo con los requisitos RF10
(API RESTful para integración ERP) y RNF08 (principios SOLID y cobertura
de pruebas).

\section{Capítulo 2: Planificación y
Análisis}\label{capuxedtulo-2-planificaciuxf3n-y-anuxe1lisis}

0.4 Análisis

La plataforma SaaS actúa como intermediaria inteligente entre tres
actores principales:

\begin{quote}
Empresas Proveedoras (B2B): Registran sus productos, configuran reglas
de venta y exámenes de certificación para vendedores.

Vendedores Freelance: Actúan como intermediarios físicos, visitan
tiendas, levantan pedidos y obtienen comisiones por cada entrega
confirmada.

Tiendas / Negocios destino: Beneficiarios finales que reciben el
producto sin necesidad de interactuar directamente con la plataforma
digital.
\end{quote}

El sistema también contempla un Administrador de Plataforma encargado de
la gestión global de usuarios, planes de suscripción y configuración del
sistema.

0.5 Requisitos Funcionales

Se identificaron 14 requisitos funcionales agrupados por módulo:

{\def\LTcaptype{none} % do not increment counter
\begin{longtable}[]{@{}
  >{\raggedright\arraybackslash}p{(\linewidth - 4\tabcolsep) * \real{0.1034}}
  >{\raggedright\arraybackslash}p{(\linewidth - 4\tabcolsep) * \real{0.2495}}
  >{\raggedright\arraybackslash}p{(\linewidth - 4\tabcolsep) * \real{0.5197}}@{}}
\toprule\noalign{}
\begin{minipage}[b]{\linewidth}\raggedright
\begin{quote}
ID
\end{quote}
\end{minipage} & \begin{minipage}[b]{\linewidth}\raggedright
Nombre
\end{minipage} & \begin{minipage}[b]{\linewidth}\raggedright
Descripción
\end{minipage} \\
\midrule\noalign{}
\endhead
\bottomrule\noalign{}
\endlastfoot
\begin{minipage}[t]{\linewidth}\raggedright
\begin{quote}
RF01
\end{quote}
\end{minipage} & Registro de Empresa Proveedora & El sistema permite el
registro de empresas con validación de datos legales (RUC/SRI) y
aprobación manual antes de activar la cuenta. \\
\begin{minipage}[t]{\linewidth}\raggedright
\begin{quote}
RF02
\end{quote}
\end{minipage} & Gestión de Catálogo de Productos & Las empresas
proveedoras pueden crear, editar y desactivar productos con precio,
stock, categoría, imagen y límites mínimos de compra. \\
\begin{minipage}[t]{\linewidth}\raggedright
\begin{quote}
RF03
\end{quote}
\end{minipage} & Registro de Vendedor Freelance & El sistema permite el
registro de vendedores con datos personales, biometría opcional y estado
de verificación pendiente. \\
\begin{minipage}[t]{\linewidth}\raggedright
\begin{quote}
RF04
\end{quote}
\end{minipage} & Sistema de

Certificaciones y

Exámenes & Las empresas pueden crear exámenes de certificación; los
vendedores deben aprobarlos para desbloquear categorías de productos
restringidos. \\
\begin{minipage}[t]{\linewidth}\raggedright
\begin{quote}
RF05
\end{quote}
\end{minipage} & Visualización de Comisiones & Antes de confirmar un
pedido, el vendedor visualiza el desglose detallado de su comisión
estimada (porcentaje, monto y condiciones). \\
\begin{minipage}[t]{\linewidth}\raggedright
\begin{quote}
RF06
\end{quote}
\end{minipage} & Gestión de Pedidos con Anticipo & El sistema soporta el
flujo de pedido con pago de anticipo (50\% o 100\%) obligatorio,
bloqueando la generación del pedido si no se cumple. \\
\begin{minipage}[t]{\linewidth}\raggedright
\begin{quote}
RF07
\end{quote}
\end{minipage} & Flujo de Pago Directo Empresa-Comprador & El dinero del
cliente final se transfiere directamente a la empresa proveedora; la
plataforma no retiene fondos del cliente. \\
\end{longtable}
}

Continua en la siguiente página...

(Continua de la página anterior)

{\def\LTcaptype{none} % do not increment counter
\begin{longtable}[]{@{}
  >{\raggedright\arraybackslash}p{(\linewidth - 4\tabcolsep) * \real{0.1034}}
  >{\raggedright\arraybackslash}p{(\linewidth - 4\tabcolsep) * \real{0.2495}}
  >{\raggedright\arraybackslash}p{(\linewidth - 4\tabcolsep) * \real{0.5197}}@{}}
\toprule\noalign{}
\begin{minipage}[b]{\linewidth}\raggedright
\begin{quote}
ID
\end{quote}
\end{minipage} & \begin{minipage}[b]{\linewidth}\raggedright
Nombre
\end{minipage} & \begin{minipage}[b]{\linewidth}\raggedright
Descripción
\end{minipage} \\
\midrule\noalign{}
\endhead
\bottomrule\noalign{}
\endlastfoot
\begin{minipage}[t]{\linewidth}\raggedright
\begin{quote}
RF08
\end{quote}
\end{minipage} & Liberación de

Comisión por Entrega

Confirmada & La comisión del vendedor se libera únicamente cuando la
empresa proveedora confirma la entrega exitosa del pedido en la tienda
destino. \\
\begin{minipage}[t]{\linewidth}\raggedright
\begin{quote}
RF09
\end{quote}
\end{minipage} & Sistema de Calificación

Bidireccional & Vendedores y empresas pueden calificarse mutuamente (1 5
estrellas con comentario) tras la confirmación de entrega de cada
pedido. \\
\begin{minipage}[t]{\linewidth}\raggedright
\begin{quote}
RF10
\end{quote}
\end{minipage} & API RESTful para Integración ERP & El sistema expone
endpoints de API RESTful documentados (OpenAPI/Swagger) para integrar
inventario, pedidos y estados de entrega con ERP externos. \\
\begin{minipage}[t]{\linewidth}\raggedright
\begin{quote}
RF11
\end{quote}
\end{minipage} & Módulo Nativo de

Inventario & Para empresas sin ERP, se ofrece un módulo nativo de
gestión de inventario con alertas de stock bajo. \\
\begin{minipage}[t]{\linewidth}\raggedright
\begin{quote}
RF12
\end{quote}
\end{minipage} & Panel de Control por Rol & Cada tipo de usuario
(Empresa, Vendedor, Administrador) tiene un panel de control
diferenciado con las funcionalidades correspondientes a su rol. \\
\begin{minipage}[t]{\linewidth}\raggedright
\begin{quote}
RF13
\end{quote}
\end{minipage} & Gestión de

Suscripciones y

Facturación & El sistema gestiona planes de suscripción para empresas o
calcula y cobra el fee por transacción (1-2\%) de manera automática. \\
\begin{minipage}[t]{\linewidth}\raggedright
\begin{quote}
RF14
\end{quote}
\end{minipage} & Notificaciones y Alertas & El sistema envía
notificaciones

(correo/plataforma) para eventos clave: pedido creado, anticipo
recibido, entrega confirmada, comisión liberada y examen disponible. \\
\end{longtable}
}

0.6 Requisitos No Funcionales

Se identificaron 12 requisitos no funcionales:

{\def\LTcaptype{none} % do not increment counter
\begin{longtable}[]{@{}
  >{\raggedright\arraybackslash}p{(\linewidth - 4\tabcolsep) * \real{0.0119}}
  >{\raggedright\arraybackslash}p{(\linewidth - 4\tabcolsep) * \real{0.0917}}
  >{\raggedright\arraybackslash}p{(\linewidth - 4\tabcolsep) * \real{0.2217}}@{}}
\toprule\noalign{}
\multicolumn{2}{@{}>{\raggedright\arraybackslash}p{(\linewidth - 4\tabcolsep) * \real{0.1036} + 2\tabcolsep}}{%
\begin{minipage}[b]{\linewidth}\raggedright
\begin{quote}
ID
\end{quote}
\end{minipage}} & \begin{minipage}[b]{\linewidth}\raggedright
Categoría
\end{minipage} \\
\midrule\noalign{}
\endhead
\bottomrule\noalign{}
\endlastfoot
RNF01 & Rendimiento & Páginas principales cargan en menos de 3 segundos
bajo condiciones normales. Consultas de catálogo e inventario responden
en menos de 2 segundos. \\
RNF02 & Disponibilidad & Disponibilidad mínima del 99,5\% mensual
(máximo ≈3,6 horas de inactividad/mes). \\
RNF03 & Escalabilidad & La arquitectura de microservicios permite
escalabilidad horizontal independiente de los módulos de Inventario,
Pagos y Certificaciones ante picos de demanda. \\
\end{longtable}
}

(Continua de la página anterior)

{\def\LTcaptype{none} % do not increment counter
\begin{longtable}[]{@{}
  >{\raggedright\arraybackslash}p{(\linewidth - 4\tabcolsep) * \real{0.1034}}
  >{\raggedright\arraybackslash}p{(\linewidth - 4\tabcolsep) * \real{0.2212}}
  >{\raggedright\arraybackslash}p{(\linewidth - 4\tabcolsep) * \real{0.5366}}@{}}
\toprule\noalign{}
\begin{minipage}[b]{\linewidth}\raggedright
\begin{quote}
ID
\end{quote}
\end{minipage} & \begin{minipage}[b]{\linewidth}\raggedright
Categoría
\end{minipage} & \begin{minipage}[b]{\linewidth}\raggedright
Descripción
\end{minipage} \\
\midrule\noalign{}
\endhead
\bottomrule\noalign{}
\endlastfoot
\begin{minipage}[t]{\linewidth}\raggedright
\begin{quote}
RNF04
\end{quote}
\end{minipage} & Seguridad

Autenticación & Autenticación con tokens JWT y soporte de 2FA para
cuentas de empresa. Las sesiones expiran tras inactividad. \\
\begin{minipage}[t]{\linewidth}\raggedright
\begin{quote}
RNF05
\end{quote}
\end{minipage} & Seguridad

Datos Financieros & Todos los datos de pago se transmiten bajo

HTTPS/TLS 1.2+. No se almacenan datos de tarjeta en crudo; se delega a
pasarela certificada PCI-DSS. \\
\begin{minipage}[t]{\linewidth}\raggedright
\begin{quote}
RNF06
\end{quote}
\end{minipage} & Usabilidad

(Vendedores) & Interfaz Mobile- rst. Botones mínimo 48×48px, fuentes no
menores a 18px, flujos con máximo 3 pasos por acción principal. \\
\begin{minipage}[t]{\linewidth}\raggedright
\begin{quote}
RNF07
\end{quote}
\end{minipage} & Usabilidad (Portal Empresarial) & Portal B2B optimizado
para escritorio con barra lateral ja, tablas paginadas y exportación de
datos (CSV/Excel). \\
\begin{minipage}[t]{\linewidth}\raggedright
\begin{quote}
RNF08
\end{quote}
\end{minipage} & Mantenibilidad & El c digo sigue principios SOLID con
cobertura de pruebas unitarias mínima del 70\% en módulos críticos
(Pagos, Certificaciones, Pedidos). \\
\begin{minipage}[t]{\linewidth}\raggedright
\begin{quote}
RNF09
\end{quote}
\end{minipage} & Cumplimiento Normativo & El flujo de dinero está
diseñado para cumplir regulaciones antilavado (AML) en Latinoamérica,
garantizando que los fondos van directamente empresa-comprador. \\
\begin{minipage}[t]{\linewidth}\raggedright
\begin{quote}
RNF10
\end{quote}
\end{minipage} & Compatibilidad & Compatible con Chrome, Firefox, Safari
y Edge (últimas 2 versiones). Vista de vendedor responsive en
dispositivos móviles. \\
\begin{minipage}[t]{\linewidth}\raggedright
\begin{quote}
RNF11
\end{quote}
\end{minipage} & Contabilidad de Transacciones & La base de datos
PostgreSQL garantiza propiedades ACID en todas las transacciones
financieras, evitando registros inconsistentes o duplicados. \\
\begin{minipage}[t]{\linewidth}\raggedright
\begin{quote}
RNF12
\end{quote}
\end{minipage} & Trazabilidad y Auditor a & Todas las acciones
financieras quedan registradas en un log de auditor a inmutable con
timestamp, usuario y estado. \\
\end{longtable}
}

0.7 Historias de Usuario / Casos de Uso

\subsection{0.7.1 Actores del Sistema}\label{actores-del-sistema}

\begin{quote}
Administrador de Plataforma: Gestiona usuarios, empresas, planes de
suscripción y configuración global del sistema.

Empresa Proveedora: Gestiona su catálogo, con gura exámenes de
certificación, con rma despachos y califica vendedores.

Vendedor Freelance: Se registra, rinde exámenes, genera pedidos en
nombre de tiendas, paga anticipos y recibe comisiones.
\end{quote}

\subsection{0.7.2 Casos de Uso Módulo de
Pedidos}\label{casos-de-uso-muxf3dulo-de-pedidos}

{\def\LTcaptype{none} % do not increment counter
\begin{longtable}[]{@{}
  >{\raggedright\arraybackslash}p{(\linewidth - 6\tabcolsep) * \real{0.1204}}
  >{\raggedright\arraybackslash}p{(\linewidth - 6\tabcolsep) * \real{0.1933}}
  >{\raggedright\arraybackslash}p{(\linewidth - 6\tabcolsep) * \real{0.2498}}
  >{\raggedright\arraybackslash}p{(\linewidth - 6\tabcolsep) * \real{0.3112}}@{}}
\toprule\noalign{}
\begin{minipage}[b]{\linewidth}\raggedright
\begin{quote}
ID
\end{quote}
\end{minipage} & \begin{minipage}[b]{\linewidth}\raggedright
Caso de Uso
\end{minipage} & \begin{minipage}[b]{\linewidth}\raggedright
Actor Principal
\end{minipage} & \begin{minipage}[b]{\linewidth}\raggedright
Descripción
\end{minipage} \\
\midrule\noalign{}
\endhead
\bottomrule\noalign{}
\endlastfoot
\begin{minipage}[t]{\linewidth}\raggedright
\begin{quote}
CU01
\end{quote}
\end{minipage} & Crear Pedido & Vendedor Freelance & Selecciona
productos del catálogo desbloqueado, especifica cantidades y tienda
destino. El sistema calcula el anticipo y la comisión estimada. \\
\begin{minipage}[t]{\linewidth}\raggedright
\begin{quote}
CU02
\end{quote}
\end{minipage} & Pagar Anticipo & Vendedor Freelance & Realiza el pago
del anticipo (50\% o 100\%) mediante pasarela. El pedido pasa a En
preparación solo si el pago es con rmado. \\
\begin{minipage}[t]{\linewidth}\raggedright
\begin{quote}
CU03
\end{quote}
\end{minipage} & Confirmar Despacho & Empresa Proveedora & La empresa
con rma la entrega exitosa en la tienda destino. El sistema libera
automáticamente la comisión al vendedor. \\
\begin{minipage}[t]{\linewidth}\raggedright
\begin{quote}
CU04
\end{quote}
\end{minipage} & Cali car & Vendedor / Empresa & Tras la confirmación,
ambas partes pueden calificarse mutuamente (1 5 estrellas con comentario
opcional). \\
\begin{minipage}[t]{\linewidth}\raggedright
\begin{quote}
CU05
\end{quote}
\end{minipage} & Consultar

Comisiones & Vendedor Freelance & El vendedor visualiza el historial de
comisiones: pendientes, liberadas y proyectadas por pedidos activos. \\
\end{longtable}
}

\subsection{0.7.3 Historias de Usuario
Representativas}\label{historias-de-usuario-representativas}

\begin{quote}
HU-01 (Vendedor): Como vendedor freelance, quiero ver el desglose de mi
comisión antes de confirmar un pedido, para decidir si el margen
justifica la visita a la tienda.

HU-02 (Empresa): Como empresa proveedora, quiero que la comisión del
vendedor se libere automáticamente solo cuando yo con rme la entrega,
para tener control sobre el pago y evitar fraudes.

HU-03 (Vendedor): Como vendedor freelance, quiero rendir un examen de
certificación para desbloquear categorías de productos especiales y
aumentar mi catálogo disponible.

HU-04 (Empresa): Como empresa proveedora, quiero integrar mi ERP con la
plataforma v a API REST para sincronizar inventario y pedidos en tiempo
real.

HU-05 (Admin): Como administrador, quiero revisar pedidos en disputa y
tomar decisiones de resolución para garantizar la confianza en la
plataforma.
\end{quote}

0.8 Sprint

La planificación de sprints se distribuye en cuatro ciclos quincenales
del primer bimestre:

{\def\LTcaptype{none} % do not increment counter
\begin{longtable}[]{@{}
  >{\raggedright\arraybackslash}p{(\linewidth - 4\tabcolsep) * \real{0.1203}}
  >{\raggedright\arraybackslash}p{(\linewidth - 4\tabcolsep) * \real{0.2495}}
  >{\raggedright\arraybackslash}p{(\linewidth - 4\tabcolsep) * \real{0.5047}}@{}}
\toprule\noalign{}
\begin{minipage}[b]{\linewidth}\raggedright
\begin{quote}
Sprint
\end{quote}
\end{minipage} & \begin{minipage}[b]{\linewidth}\raggedright
Objetivo
\end{minipage} & \begin{minipage}[b]{\linewidth}\raggedright
Entregables Principales
\end{minipage} \\
\midrule\noalign{}
\endhead
\bottomrule\noalign{}
\endlastfoot
\begin{minipage}[t]{\linewidth}\raggedright
\begin{quote}
Sprint 1
\end{quote}
\end{minipage} & Configuración base y autenticación & Proyecto Django
configurado, modelos de usuario/empresa/vendedor, JWT auth, panel de
admin básico. \\
\begin{minipage}[t]{\linewidth}\raggedright
\begin{quote}
Sprint 2
\end{quote}
\end{minipage} & Catálogo e inventario & CRUD de productos (RF02),
módulo nativo de inventario (RF11), API endpoints para ERP (RF10). \\
\begin{minipage}[t]{\linewidth}\raggedright
\begin{quote}
Sprint 3
\end{quote}
\end{minipage} & Pedidos y pagos & Flujo completo de pedido (RF06),
integración pasarela de pago, liberación de comisión (RF08),
notificaciones (RF14). \\
\begin{minipage}[t]{\linewidth}\raggedright
\begin{quote}
Sprint 4
\end{quote}
\end{minipage} & Certificaciones, calificaciones y panel de control &
Sistema de exámenes (RF04), calificación bidireccional (RF09),
dashboards por rol (RF12), pruebas y documentación. \\
\end{longtable}
}

Table 5: Planificación de sprints Primer Bimestre

0.9 Diagramas: Clases y BD (Diseño Lógico)

\subsection{0.9.1 Descripción del Diagrama de
Clases}\label{descripciuxf3n-del-diagrama-de-clases}

El diagrama de clases del sistema re eja las entidades principales y sus
relaciones. Las entidades clave son:

\begin{quote}
Usuario (base): id, email, password\_hash, rol (admin / empresa /
vendedor), fecha\_creacion, activo.

EmpresaProveedora: hereda de Usuario. Agrega ruc, razon\_social,
estado\_aprobacion, plan\_suscripcion.

VendedorFreelance: hereda de Usuario. Agrega cedula,
estado\_verificacion, saldo\_comision.

Producto: id, empresa\_id (FK), nombre, precio, stock, categoria,
activo, minimo\_compra.

Pedido: id, vendedor\_id (FK), estado (BORRADOR, PENDIENTE\_PAGO,
ANTICIPO\_RECIBIDO, EN\_PREPARACION, ENTREGADO, COMPLETADO, CANCELADO),
anticipo\_requerido, comision\_estimada.

ItemPedido: pedido\_id (FK), producto\_id (FK), cantidad,
precio\_unitario.

Pago: id, pedido\_id (FK), monto, tipo (anticipo/saldo), estado,
referencia\_pasarela.

Certi cacion: id, empresa\_id (FK), nombre, categoria\_producto.

ResultadoExamen: vendedor\_id (FK), certificacion\_id (FK), puntaje,
aprobado, fecha.

Cali cacion: pedido\_id (FK), calificador\_id, calificado\_id, estrellas
(1 5), comentario.
\end{quote}

\subsection{0.9.2 Preguntas Gu a para el Diagrama de Clases y
BD}\label{preguntas-gu-a-para-el-diagrama-de-clases-y-bd}

Las siguientes cinco preguntas orientan el diseño detallado del diagrama
de clases y el modelo relacional, y serÆn respondidas formalmente al
elaborar los diagramas en Figma y en el modelado l gico:

P1. C mo se modela la herencia de roles de usuario (Admin, Empresa,
Vendedor) en la base de datos relacional de PostgreSQL, usando Django
ORM: tabla œnica con discriminador, tablas separadas con OneToOneField o
MTI (Multi-Table Inheritance)?

\begin{quote}
AnÆlisis esperado: Evaluar la estrategia ptima en tØrminos de integridad
referencial, complejidad de consultas y compatibilidad con el sistema de
autenticaci n de Django.
\end{quote}

P2. QuØ atributos y restricciones de integridad (constraints) son
necesarios en la entidad Pedido para garantizar que un pedido no pueda
avanzar de estado sin el anticipo con rmado, evitando inconsistencias de
datos?

\begin{quote}
AnÆlisis esperado: De nir check constraints, transiciones de estado
vÆlidas y el rol de las propiedades ACID de PostgreSQL en este ujo.
\end{quote}

P3. Cómo se relacionan las entidades Certificacion, Producto y
VendedorFreelance para implementar el desbloqueo de categor as de
productos restringidos tras aprobar un examen?

\begin{quote}
AnÆlisis esperado: Diseæar la tabla intermedia o relaci n many-to-many
que permita consultar e cientemente quØ productos puede vender un
vendedor segœn sus certi caciones aprobadas.
\end{quote}

P4. Cómo se garantiza la trazabilidad nanciera inmutable de los eventos
de pago, anticipo y comisi n, cumpliendo con RNF12 (log de auditor a),
sin impactar el rendimiento de las consultas transaccionales?

\begin{quote}
AnÆlisis esperado: Evaluar el uso de una tabla de auditor a separada con
triggers o seæales de Django, versus el uso de un log append-only, y su
impacto en el diseæo del diagrama de clases.
\end{quote}

P5. QuØ estructura de datos y relaciones se necesitan en el m dulo de
Calificacion para soportar el sistema bidireccional (vendedor cali ca
empresa y empresa cali ca vendedor) sin duplicar entidades ni generar
ambig edad en las referencias de calificador\_id y calificado\_id?

\begin{quote}
AnÆlisis esperado: Analizar el uso de Generic Relations de Django,
claves forÆneas polim r cas o un diseæo con campo tipo\_calificacion
para distinguir la direcci n de la cali caci n.

\includegraphics[width=2.92708in,height=2.35976in]{./media/image1.png}\includegraphics[width=6.96875in,height=3.72639in]{./media/image2.png}\includegraphics[width=6.28403in,height=4.375in]{./media/image3.png}
\end{quote}

0.10 Diagrama de Componentes

La arquitectura de componentes define los m dulos del sistema y sus
interfaces de comunicaci n:

{\def\LTcaptype{none} % do not increment counter
\begin{longtable}[]{@{}
  >{\raggedright\arraybackslash}p{(\linewidth - 4\tabcolsep) * \real{0.2332}}
  >{\raggedright\arraybackslash}p{(\linewidth - 4\tabcolsep) * \real{0.2777}}
  >{\raggedright\arraybackslash}p{(\linewidth - 4\tabcolsep) * \real{0.3637}}@{}}
\toprule\noalign{}
\begin{minipage}[b]{\linewidth}\raggedright
\begin{quote}
Componente
\end{quote}
\end{minipage} & \begin{minipage}[b]{\linewidth}\raggedright
Responsabilidad
\end{minipage} & \begin{minipage}[b]{\linewidth}\raggedright
Tecnolog a
\end{minipage} \\
\midrule\noalign{}
\endhead
\bottomrule\noalign{}
\endlastfoot
\begin{minipage}[t]{\linewidth}\raggedright
\begin{quote}
core-auth
\end{quote}
\end{minipage} & Autenticación, JWT, roles y permisos & Django +
SimpleJWT \\
\begin{minipage}[t]{\linewidth}\raggedright
\begin{quote}
inventario
\end{quote}
\end{minipage} & Catalogo de productos, stock, alertas & Django app +
PostgreSQL \\
\begin{minipage}[t]{\linewidth}\raggedright
\begin{quote}
certificaciones
\end{quote}
\end{minipage} & Exámenes, evaluaciones, desbloqueo de catálogo & Django
app + PostgreSQL \\
\begin{minipage}[t]{\linewidth}\raggedright
\begin{quote}
pedidos
\end{quote}
\end{minipage} & Ciclo de vida del pedido, estados, validaciones &
Django app + PostgreSQL \\
\begin{minipage}[t]{\linewidth}\raggedright
\begin{quote}
pagos
\end{quote}
\end{minipage} & Anticipos, comisiones, pasarela de pagos & Django app +
Stripe/PayPhone \\
\begin{minipage}[t]{\linewidth}\raggedright
\begin{quote}
notificaciones
\end{quote}
\end{minipage} & Emails y alertas en plataforma & Celery + Redis +
SendGrid \\
\begin{minipage}[t]{\linewidth}\raggedright
\begin{quote}
api-gateway
\end{quote}
\end{minipage} & Enrutamiento, rate limiting, TLS & Nginx \\
\begin{minipage}[t]{\linewidth}\raggedright
\begin{quote}
spa-frontend
\end{quote}
\end{minipage} & Interfaz web reactiva con vistas por rol & React 18 +
Tailwind CSS \\
\end{longtable}
}

Table 6: Componentes del sistema y tecnolog as asociadas

\includegraphics[width=5.92791in,height=3.7401in]{./media/image4.png}

\includegraphics[width=5.92791in,height=3.45882in]{./media/image5.png}

0.11 Arquitectura Lógica y Física (Propuesta)

\subsection{0.11.1 Arquitectura Lógica}\label{arquitectura-luxf3gica}

La arquitectura lógica organiza el sistema en tres capas bien definidas:

\begin{enumerate}
\def\labelenumi{\arabic{enumi}.}
\item
  Capa de Presentación (Cliente): SPA en React con vistas diferenciadas
  por rol (Administrador, Empresa, Vendedor). Diseæo Mobile- rst para la
  vista de vendedor.
\item
  Capa de Negocio (API): Conjunto de aplicaciones Django expuestas v a
  Django REST Framework. Cada app Django corresponde a un dominio
  funcional (pedidos, pagos, inventario, etc.). La API Gateway (Nginx)
  centraliza el enrutamiento, el rate limiting y la terminación de TLS.
\item
  Capa de Datos: PostgreSQL 16 con esquema relacional normalizado.
  Propiedades ACID garantizadas para todas las transacciones
  financieras. CachØ con Redis para sesiones y cola de tareas asíncronas
  (Celery).
\end{enumerate}

\subsection{0.11.2 Arquitectura Física (Propuesta de
Despliegue)}\label{arquitectura-fuxedsica-propuesta-de-despliegue}

\begin{quote}
Servidor de aplicación: Contenedores Docker con Gunicorn (Django WSGI)
detrás de Nginx.
\end{quote}

Base de datos: Instancia PostgreSQL con réplica de lectura para
reportes.

\begin{quote}
Broker de tareas: Redis para Celery (notificaciones asíncronas,
liberación de comisiones).

CDN / Estáticos: Archivos estáticos y media servidos mediante CDN.

CI/CD: GitHub Actions con pipeline de pruebas automáticas y despliegue
en entorno de staging.
\end{quote}

\subsection{0.11.3 Flujo de Proceso del
Pedido}\label{flujo-de-proceso-del-pedido}

{\def\LTcaptype{none} % do not increment counter
\begin{longtable}[]{@{}
  >{\centering\arraybackslash}p{(\linewidth - 4\tabcolsep) * \real{0.1372}}
  >{\raggedright\arraybackslash}p{(\linewidth - 4\tabcolsep) * \real{0.2212}}
  >{\raggedright\arraybackslash}p{(\linewidth - 4\tabcolsep) * \real{0.5160}}@{}}
\toprule\noalign{}
\begin{minipage}[b]{\linewidth}\raggedright
\begin{quote}
Paso
\end{quote}
\end{minipage} & \begin{minipage}[b]{\linewidth}\raggedright
Estado
\end{minipage} & \begin{minipage}[b]{\linewidth}\raggedright
Descripción
\end{minipage} \\
\midrule\noalign{}
\endhead
\bottomrule\noalign{}
\endlastfoot
1 & BORRADOR & El vendedor selecciona productos, tienda destino y
confirma el carrito. \\
2 & PENDIENTE DE PAGO & Sistema calcula anticipo obligatorio. Vendedor
recibe instrucciones de pago. \\
3 & ANTICIPO

RECIBIDO & Pasarela con rma el anticipo. Sistema notifica a la
empresa. \\
4 & EN

PREPARACIÓN & Empresa prepara y despacha el pedido hacia la tienda
destino. \\
5 & ENTREGADO & Empresa con rma la entrega. Sistema libera
automáticamente la comisión. \\
6 & COMPLETADO & Ambas partes califican. Pedido archivado en
historial. \\
\multicolumn{2}{@{}>{\centering\arraybackslash}p{(\linewidth - 4\tabcolsep) * \real{0.3585} + 2\tabcolsep}}{%
\begin{minipage}[t]{\linewidth}\centering
\begin{quote}
CANCELADO /

EN DISPUTA
\end{quote}
\end{minipage}} & Si el pago falla o la entrega no se con rma en el
plazo, el pedido entra en disputa con revisión del administrador. \\
\end{longtable}
}

Table 7: Estados del ciclo de vida del pedido

0.12 Prototipo de Pantallas No Funcional

El prototipo de pantallas no funcional será desarrollado en Figma y
constituirá la base visual para el posterior diseño del diagrama de
clases y la arquitectura de la base de datos. El prototipo contemplará
las siguientes vistas:

\begin{quote}
Pantalla de inicio de sesión (roles diferenciados).

Dashboard del Administrador: gestión de empresas, vendedores y planes.

Dashboard de la Empresa Proveedora: catálogo de productos, pedidos
activos, confirmación de despachos y calificaciones.

Dashboard del Vendedor Freelance (Mobile- rst): catálogo disponible,
creación de pedido, vista de comisiones y exámenes de certificación.

Flujo de creación de pedido: selección de productos cálculo de anticipo
y comisión pago confirmación.

Módulo de exámenes: listado de certificaciones disponibles y pantalla de
evaluación.
\end{quote}

Nota: El enlace al prototipo en Figma será adjuntado en el repositorio
de GitHub dentro de la carpeta PRESENTACION-1BIM, una vez finalizado el
diseño de pantallas.

\section{Referencias}\label{referencias}

\begin{quote}
Django Software Foundation. (2024). Django 5.x Documentation.
\href{https://docs.djangoproject.com/}{https:// docs.djangoproject.com/}

Encode OSS Ltd. (2024). Django REST Framework Documentation.
\href{https://www.django-rest-framework.org/}{https:
//www.django-rest-framework.org/} PostgreSQL Global Development Group.
(2024). PostgreSQL 16 Documentation.
\url{https://www.postgresql.org/docs/}

Parque Científico y Tecnológico UTPL. (2026). Análisis de Solución del
Reto Plataforma SaaS de Distribución Freelance. Versión 1.0.
\end{quote}

\end{document}
